\documentclass{article}

\usepackage{amsmath}
\usepackage{amssymb}

\newcommand{\ZZ}{\mathbb{Z}}
\newcommand{\RR}{\mathbb{R}}
\newcommand{\CC}{\mathbb{C}}

\begin{document}

{\bf Worksheet \#8; date: 09/25/2018}

{\bf MATH 55 Discrete Mathematics}

\begin{enumerate}
\item {\em (Rosen 4.4.6c)} Find an inverse of $a$ modulo $m$ for the following pairs of relatively prime integers.
\[
a = 144, \quad m = 233.
\]

\item {\em (Rosen 4.4.7)} Show that if $a$ and $m$ are relatively prime positive integers, then the inverse of $a$ modulo $m$ is unique modulo $m$. {\em Hint:} Assume that there are two solutions $b$ and $c$ of the congruence $ax \equiv 1 \pmod{m}$, and show that $b \equiv c \pmod{m}$.

\item {\em (Rosen 4.4.12b)} Solve the following congruences using the modular inverses found in Question 1:
\[
144x \equiv 4 \pmod{233}
\]

\item {\em (Rosen 4.4.21, 24)} Find all solutions to the system of congruences $x \equiv 1 \pmod{2}$, $x \equiv 2 \pmod{3}$, $x \equiv 3 \pmod{5}$, and $x \equiv 4 \pmod{11}$ with two methods:
\begin{enumerate}
\item The construction in the proof of the Chinese remainder theorem
\item Back substitution.
\end{enumerate}

\item {\em (Challeging)} This questions put all techniques above together. Try this now, or after the lecture!
\[
\left\{
\begin{array}{rcl}
9x + 1 &\equiv& 0 \pmod{2} \\
5x &\equiv& 1 \pmod{3} \\
x &\equiv& 4 \pmod{7}
\end{array}
\right.
\]
\end{enumerate}

\end{document}